\scanpage{15}

% Detached mark at whole-source approximately (200,103), below-left of J;
% its textual status is unresolved, not its physical presence.
\makebox[0pt][l]{\raisebox{-1.3ex}{\hspace{-0.25em}\unsure{.|\phantom{.}}}}Je te signale que j'ai \unsure{réfléchi} aux Gr-champs sur $X$. Si $G$ est un Groupe sur $X$; $N$ un $G$-Module les Gr-champs sur $X$ ``épinglés par $G,N$'' forment a priori une $2$-catégori\unsure{\phantom{e}|e} et même une $2$-catégorie de Picard stricte, grâce à l'opération \unsure{\.{e}vident|évident} à la Baer. On trouve que le complexe \inserted{(de cochaines)} tronqué à l'échelon $2$ qui leur correspond est le tronqué
\[
\tau_{\leq 2}\bigl(\mathbb R\Gamma(B_{G}\mathbin{\mathrm{mod}}X,N)[1]\bigr)
\]
(\uline{NB} la cohomologie de $\mathbb R\Gamma(B_{G}\mathbin{\mathrm{mod}}X,N)$ commence en degré $1$) Plus \unsure{géométriquement}, un Gr-champ sur $X$ \unsure{\.{e}|é}pi\unsure{-|n}g\unsure{\phantom{lé}|lé} \inserted{par $(G,N)$} est essentiellement ``la même chose'' qu'une $2$-gerbe sur $B_{G}$, liée par $N$, et \deleted{\unsure{épi|epi|pri}}munie d'une trivialisati\unsure{\phantom{on}|on} au dessus de $X\approx B_{c}=(B_{G})_{/P}$ (où $P$ est le \makebox[0pt][l]{\raisebox{.9ex}{\hspace{-.15em}\deleted{\unsure{s|$\ast$}}}}objet de $B_{G}$ ``torseur universel sous $G$''). Ces $2$-gerbes forment en fait une $3$-catégorie
% Detached point below de, approximately (175,1630); the underset is positional.
$\underset{\unsure{\cdot|\phantom{\cdot}}}{\text{de}}$ Picard a priori, mais il se
% archive sheet number: 14
